\section{Intermediate Upper Bounds}\label{sec:intermediate}

The main upper bound in \Cref{sec:main-upper} refines a simpler
odd-prime-prefix method.  We present the two intermediate versions because
their proofs isolate the compression and Bonferroni ideas used later.

\subsection{Odd-part compression}

The following elementary lemma is used throughout the upper-bound arguments.

\begin{lemma}[Odd-part injection]\label{lem:odd-part-injection}
Let $A$ be a divisibility antichain.  The map
\[
  \varphi(x):=\frac{x}{2^{v_2(x)}}
\]
is injective on $A$.
\end{lemma}

\begin{proof}
If $\varphi(x)=\varphi(y)=m$, then $x=2^a m$ and $y=2^b m$.  If $a<b$, then
$x\mid y$; if $b<a$, then $y\mid x$.  Since $A$ is an antichain, both cases
force $x=y$.
\end{proof}

\subsection{The \texorpdfstring{$13/36$}{13/36} bound}

\begin{theorem}\label{thm:thirteen-thirtysix}
Shortener has a strategy forcing
\[
  \L(n)\le\left(\frac{13}{36}+o(1)\right)n.
\]
\end{theorem}

\begin{proof}
Let
\[
  k:=\left\lfloor\frac{\sqrt n}{\log n}\right\rfloor.
\]
For her first $k$ turns, Shortener plays the smallest legal odd prime; after
that she plays arbitrarily.  Let
\[
  q_1<q_2<\cdots<q_k
\]
be the odd primes she obtains.

We first bound these primes.  Fix $\varepsilon>0$.  If, before Shortener's
$j$th prefix move, every odd prime up to
\[
  X=(3/2+\varepsilon)j\log n
\]
were blocked, then the total odd-prime Chebyshev mass up to $X$ would be
covered by the earlier Shortener primes together with the odd prime divisors of
Prolonger's first $j$ moves.  Each Prolonger move contributes at most
$\log n$ of prime-divisor log-mass.  By induction on $j$, the previous
Shortener primes contribute at most
\[
  \sum_{i<j}\log q_i\le \left(\frac12+o(1)\right)j\log n,
\]
because $j\le\sqrt n/\log n$.  Thus the blocked log-mass is at most
$(3/2+o(1))j\log n$, while the prime number theorem gives
\[
  \vartheta_{\rm odd}(X)=(3/2+\varepsilon+o(1))j\log n,
\]
a contradiction for large $n$.  Hence
\[
  q_j\le(3/2+\varepsilon)j\log n
\]
uniformly for $j\le k$, and therefore
\[
  \sum_{j=1}^k\frac1{q_j}\ge\frac13-o(1).
  \tag{6.1}\label{eq:one-third-mass}
\]

Let $A$ be the final antichain.  All elements other than the prefix primes have
odd part not divisible by any $q_j$.  By \Cref{lem:odd-part-injection},
\[
  |A|\le k+N_D(n),
\]
where $D=\{q_1,\ldots,q_k\}$ and $N_D(n)$ is the number of odd integers
$m\le n$ not divisible by any prime in $D$.

Choose $t$ minimal such that
\[
  s_t:=\sum_{j\le t}\frac1{q_j}\ge\frac13-o(1),
\]
and put $E=\{q_1,\ldots,q_t\}$.  Since every $q_j$ is an odd prime,
$1/q_j\le1/3$, so
\[
  s_t\in\left[\frac13-o(1),\,\frac23+o(1)\right]\subset[0,1]
\]
for large $n$.  Because $E\subseteq D$, we have $N_D(n)\le N_E(n)$.

Second-order Bonferroni on the events $q\mid m$ over odd $m\le n$ gives
\[
  N_E(n)
  \le
  \frac n2\left(1-s_t+\frac{s_t^2}{2}\right)+O(t^2).
\]
Here $t\le k$, so $O(t^2)=O(n/\log^2n)=o(n)$.  The function
$f(s)=1-s+s^2/2$ is decreasing on $[0,1]$, and therefore
\[
  f(s_t)\le f(1/3-o(1))=\frac{13}{18}+o(1).
\]
Consequently
\[
  |A|\le k+N_D(n)\le\frac{13}{36}n+o(n).
\]
\end{proof}

\subsection{The \texorpdfstring{$5/16$}{5/16} bound}

\begin{theorem}\label{thm:five-sixteen}
Shortener has a strategy forcing
\[
  \L(n)\le\left(\frac{5}{16}+o(1)\right)n.
\]
\end{theorem}

\begin{proof}
Fix $A_\star>2$ and set
\[
  k:=\left\lfloor\frac{n}{2A_\star\log n}\right\rfloor.
\]
Shortener plays the smallest legal odd prime for her first $k$ turns.  The
same Chebyshev log-mass induction, now with $j\le n/(2A_\star\log n)$, gives
\[
  q_j\le A_\star j\log n
  \qquad (1\le j\le k).
\]
Indeed, if all odd primes below $A_\star j\log n$ were blocked, Prolonger's first
$j$ moves would contribute at most $j\log n$, and the previous Shortener primes
would contribute at most $(1+o(1))j\log n$; this is less than
$A_\star j\log n$ because $A_\star>2$.

It follows that
\[
  \sum_{j=1}^k\frac1{q_j}\ge \frac1{A_\star}-o(1).
\]
As in the proof of \Cref{thm:thirteen-thirtysix}, use the odd-part injection and
truncate the prime set at the first partial sum
\[
  s_t\ge \frac1{A_\star}-o(1).
\]
Since $1/q_j\le1/3$ and $A_\star>2$, the partial sum remains in $[0,1]$.  The
second-order Bonferroni estimate gives
\[
  \L(n)
  \le
  \frac n2
  \left(1-\frac1{A_\star}+\frac{1}{2A_\star^2}\right)
  +o(n).
\]
Letting $A_\star\downarrow2$ yields
\[
  \L(n)\le \frac n2\left(1-\frac12+\frac18\right)+o(n)
  =
  \frac{5}{16}n+o(n).
\]
The pair-intersection error is $o(n)$ because only prime pairs with product at
most $n$ contribute an $O(1)$ floor error, and
\[
  \#\{p<q:\ pq\le n\}
  \ll \sum_{p\le\sqrt n}\pi(n/p)
  \ll \frac{n}{\log n}\sum_{p\le\sqrt n}\frac1p
  =
  o(n).
\]
\end{proof}
